% Mini-relatório: apresentação acadêmica baseada nas normas ABNT.
% O cabeçalho substitui a capa separada para respeitar as cinco páginas.
% gerar_relatorio.py preenche os dados e incorpora os gráficos em TikZ.
\documentclass[12pt,a4paper]{article}
\usepackage[T1]{fontenc}
\usepackage[utf8]{inputenc}
\usepackage[brazil]{babel}
\usepackage{newtxtext,newtxmath}
\usepackage[left=3cm,top=3cm,right=2cm,bottom=2cm,headheight=14pt,headsep=1cm]{geometry}
\usepackage{setspace,indentfirst,booktabs,array,tikz,titlesec,fancyhdr}
\usepackage[hyphens]{url}
\usepackage[hidelinks]{hyperref}
\urlstyle{same}
\onehalfspacing
\setlength{\parindent}{1.25cm}
\setlength{\parskip}{0pt}
\setlength{\emergencystretch}{2em}
\titleformat{\section}{\normalfont\normalsize\bfseries}{\thesection}{1em}{\MakeUppercase}
\titlespacing*{\section}{0pt}{\baselineskip}{0.5\baselineskip}
\pagestyle{fancy}
\fancyhf{}
\fancyhead[R]{\fontsize{10}{12}\selectfont\thepage}
\renewcommand{\headrulewidth}{0pt}
\newcommand{\fonte}{\par\noindent{\fontsize{10}{12}\selectfont Fonte: elaboração própria (2026).}\par}
\newcommand{\titulofigura}[2]{\noindent{\fontsize{10}{12}\selectfont Figura #1: #2}\par\smallskip}
\hypersetup{pdftitle={Análise de uma amostra da rede social do Facebook},
  pdfauthor={%%AUTORES_PDF%%}}
%%CORES%%
\begin{document}
\begin{center}
\begin{singlespace}
\textbf{UNIVERSIDADE DE BRASÍLIA}\\
Departamento de Ciência da Computação\\
Teoria e Aplicação de Grafos (Turma 01, 2026/2)\\[0.4cm]
\textbf{Análise de uma amostra da rede social do Facebook}\\[0.3cm]
%%AUTORES%%
Professor: Díbio\qquad Brasília, 2026
\end{singlespace}
\end{center}

\section{Introdução}
Analisamos 404 usuários da rede de amizades do Facebook. Calculamos as
centralidades de grau, intermediação, proximidade e autovetor e usamos Louvain
para encontrar comunidades. Os dados vêm do conjunto \textit{ego-Facebook}
do Stanford Network Analysis Project (2012), formado por dez redes-ego com
identificadores anônimos. Comparamos os resultados para entender quais
usuários têm mais contatos e quais ocupam posições importantes nos caminhos
entre os demais.

\section{Preparação dos dados}
Construímos um grafo simples, não direcionado e sem pesos com NetworkX
(NetworkX Developers, 2026). Lemos as amizades dos arquivos \texttt{.edges}
e acrescentamos as ligações de cada ego aos amigos da primeira coluna de
\texttt{.feat}. A união das dez redes resultou em 4.039 nós e 88.234 arestas.

Sorteamos 404 nós sem reposição, com a semente %%SEMENTE%%. Depois,
selecionamos as arestas com as duas extremidades entre os sorteados,
obtendo o subgrafo induzido. Mantivemos os nós que ficaram isolados para
analisar todos os 404 usuários. A Tabela 1 reúne as características da amostra.

\begin{singlespace}
\noindent{\fontsize{10}{12}\selectfont Tabela 1: Características da amostra}
\begin{center}
\fontsize{10}{12}\selectfont
\begin{tabular*}{\textwidth}{@{\extracolsep{\fill}}lr}
\toprule
\textbf{Característica} & \textbf{Valor}\\
\midrule
%%RESUMO%%
\bottomrule
\end{tabular*}
\end{center}
\fonte
\end{singlespace}

\newpage
\section{Centralidades e usuários de destaque}
Calculamos as quatro medidas com as funções do NetworkX (NetworkX Developers,
2026). O grau de centralidade corresponde ao número de vizinhos dividido por
$N-1$. A intermediação normalizada mede a participação do usuário nos caminhos
mínimos entre pares de outros nós. A proximidade considera a distância média
aos nós alcançáveis e usa a correção de Wasserman-Faust para levar em conta
essa fração da rede. O autovetor dá mais valor às ligações com usuários também
centrais. Ele é obtido de $Ax=\lambda x$, com normalização euclidiana.

\begin{singlespace}
\noindent{\fontsize{10}{12}\selectfont Tabela 2: Maior valor de cada centralidade}
\begin{center}
\fontsize{10}{12}\selectfont
\begin{tabular*}{\textwidth}{@{\extracolsep{\fill}}lrr}
\toprule
\textbf{Centralidade} & \textbf{Nó} & \textbf{Valor}\\
\midrule
%%LIDERES%%
\bottomrule
\end{tabular*}
\end{center}
\fonte
\end{singlespace}

O nó %%LIDER_GRAU%% possui %%GRAU_LIDER%% vizinhos em uma componente de
%%COMPONENTE_GRAU%% usuários e lidera o grau. Isso indica o maior alcance
inicial para uma informação. Já o nó %%LIDER_INTERMEDIACAO%%, com apenas
%%GRAU_INTERMEDIACAO%% vizinhos, lidera intermediação e proximidade, seguido
por %%SEGUNDO_INTERMEDIACAO%% nas duas medidas. Seus vizinhos estão em
%%COMUNIDADES_VIZINHAS%% comunidades diferentes: ele participa de muitos
caminhos mínimos e funciona como ponte entre grupos da componente.

No autovetor, o maior valor é do nó %%LIDER_AUTOVETOR%%, que pertence à mesma
componente de %%COMPONENTE_AUTOVETOR%% nós. Ele tem menos vizinhos que
%%LIDER_GRAU%%, mas seus vizinhos também têm valores altos de autovetor e
contribuem para seu resultado. Por isso, os maiores valores das quatro centralidades
nem sempre pertencem ao mesmo usuário.

O grafo tem %%NUM_COMPONENTES%% componentes conexas. Como há usuários sem
caminho entre si, a proximidade usa apenas os alcançáveis, com a correção
mencionada acima. O autovetor global ficou concentrado na maior componente,
onde estão os líderes das quatro medidas. Ao todo, %%AUTOVETORES_PEQUENOS%% nós tiveram
valor inferior a $10^{-6}$, %%AUTOVETORES_PEQUENOS_FORA%% deles fora dessa
componente; isso reflete a componente de cada um, não irrelevância local.
Usamos iteração de potência com tolerância $10^{-10}$ e limite de 1.000
iterações. Os isolados receberam zero nas quatro medidas.

\newpage
\section{Visualização das centralidades}
Na Figura 1, mantivemos as mesmas posições dos nós nos quatro mapas para
facilitar a comparação. Nós maiores têm maior grau. As cores vão do roxo,
para os menores valores da medida, ao amarelo, para os maiores. Os cinco
usuários com os maiores valores estão identificados ao lado de cada mapa.

\begin{singlespace}
\titulofigura{1}{Mapas das quatro centralidades e respectivos usuários de destaque}
\begin{center}
%%MAPAS%%
\end{center}
\fonte
\end{singlespace}

O maior grau é do nó %%LIDER_GRAU%%; intermediação e proximidade são
lideradas por %%LIDER_INTERMEDIACAO%%, e o autovetor, por %%LIDER_AUTOVETOR%%,
todos na mesma componente. O notebook traz os dez maiores valores de cada
medida. Os resultados de todos os usuários estão no arquivo \texttt{metricas.csv}.

\newpage
\section{Comunidades pelo algoritmo de Louvain}
Executamos Louvain com a função do NetworkX, resolução 1 e semente
%%SEMENTE%%. O algoritmo move nós entre comunidades quando há ganho de
modularidade e depois reúne as comunidades em uma rede menor. Essas duas
etapas se repetem até que o ganho seja pequeno (Blondel \textit{et al.}, 2008).
Encontramos %%NUM_COMUNIDADES%% comunidades, com modularidade de
%%MODULARIDADE%%. O mapa e os tamanhos dos grupos estão na Figura 2.

\begin{singlespace}
\titulofigura{2}{Comunidades detectadas por Louvain na amostra}
\begin{center}
%%MAPA_COMUNIDADES%%
\end{center}
\fonte
\end{singlespace}

Dos grupos encontrados, %%COMUNIDADES_UNITARIAS%% são formados por um único
nó isolado. Todos aparecem em cinza, e cada um tem sua própria comunidade.
A modularidade alta também está relacionada à separação do grafo em
componentes sem ligações entre si. Louvain pode dividir uma componente em
mais de uma comunidade. Os grupos foram calculados a partir das arestas;
não usamos os círculos de amizade informados nos arquivos \texttt{.circles}.

\newpage
\section{Considerações finais}
O nó %%LIDER_GRAU%% liderou o grau, %%LIDER_INTERMEDIACAO%% liderou
intermediação e proximidade, e %%LIDER_AUTOVETOR%% liderou o autovetor.
Para espalhar uma informação pelas amizades da amostra, %%LIDER_GRAU%% teria
mais contatos diretos e %%LIDER_INTERMEDIACAO%% ligaria grupos diferentes,
mas a mensagem ficaria restrita à componente de ambos.
Essa interpretação vem das ligações do grafo, pois não simulamos a
propagação de mensagens. A Figura 3 mostra a distribuição dos graus e o
tamanho das comunidades encontradas.

\begin{singlespace}
\titulofigura{3}{Distribuição dos graus e tamanhos das comunidades}
\begin{center}
%%DISTRIBUICOES%%
\end{center}
\fonte
\end{singlespace}

Ao selecionar 404 dos 4.039 usuários, perdemos ligações que existiam na rede
completa. Isso ajuda a explicar os nós isolados e as várias componentes
conexas. Assim, os resultados se referem aos usuários sorteados e às
amizades que permaneceram na amostra. O notebook permite repetir os cálculos,
e os arquivos CSV registram a seleção dos nós e os valores das métricas.

\section*{Referências}
\begin{singlespace}
\setlength{\parindent}{0pt}
\raggedright
BLONDEL, Vincent D.; GUILLAUME, Jean-Loup; LAMBIOTTE, Renaud; LEFEBVRE,
Etienne. Fast unfolding of communities in large networks.
\textbf{Journal of Statistical Mechanics: Theory and Experiment},
[S. l.], v. 2008, n. 10, P10008, 2008.
DOI: \url{https://doi.org/10.1088/1742-5468/2008/10/P10008}.
\par\medskip
NETWORKX DEVELOPERS. \textbf{NetworkX documentation}. Versão %%VERSAO_NETWORKX%%.
[S. l.], 2026. Disponível em: \url{https://networkx.org/documentation/stable/}.
Acesso em: 12 out. 2026.
\par\medskip
STANFORD NETWORK ANALYSIS PROJECT. \textbf{Social circles: Facebook}.
Stanford, 2012. Disponível em:
\url{https://snap.stanford.edu/data/egonets-Facebook.html}.
Acesso em: 12 out. 2026.
\end{singlespace}
\end{document}
